\documentclass[10pt, a4paper]{article}
\usepackage[T2A]{fontenc}
\usepackage[utf8]{inputenc}
\usepackage{enumitem}
\usepackage[russian]{babel}
\usepackage{amsmath, amsfonts, amssymb, amsthm}
\usepackage{geometry}
\geometry{top=1.5cm, bottom=1.5cm, left=1.5cm, right=1.5cm}

\title{\vspace{-1cm}\textbf{Конспект по дискретной математике}\vspace{-0.5cm}}
\author{}
\date{}

\newcommand{\divs}{\mathbin{\vcenter{\baselineskip=2.2pt\lineskiplimit=0pt\hbox{.}\hbox{.}\hbox{.}}}}
\newcommand{\ndivs}{\mathbin{\not\vcenter{\baselineskip=2.2pt\lineskiplimit=0pt\hbox{.}\hbox{.}\hbox{.}}}}
\def\letus{%
\mathord{\setbox0=\hbox{$\exists$}%
    \hbox{\kern 0.125\wd0%
        \vbox to \ht0{%
            \hrule width 0.75\wd0%
            \vfill%
            \hrule width 0.75\wd0}%
        \vrule height \ht0%
        \kern 0.125\wd0}%
    }%
}

\begin{document}
\maketitle
\vspace{-0.8cm}

\section*{Список используемых обозначений}
\begin{tabular}{ll}
    $]$ & --- Пусть \\
    Р/м & --- Рассмотрим \\
    $\lfloor x \rfloor$ & --- Округление числа $x$ вниз \\
    $\lceil x \rceil$ & --- Округление числа $x$ вверх \\
    $\iff$ & --- Тогда и только тогда, когда \\

\end{tabular}

\vspace{0.5cm}
\hrule
\vspace{0.5cm}

\centerline{\Large\textbf{\S \,\, Методы доказательств}}
\vspace{0.4cm}

\subsection*{Конструктивное док-во (пример)}

\textbf{Пример:} Докажите, что пифагоровых троек бесконечно много.
\[ 3^2 + 4^2 = 5^2 \]
\[ (3k)^2 + (4k)^2 = (5k)^2 \text{ --- бесконечная серия, } \forall k \in \mathbb{N} \]

\subsection*{Контрпример}

\textbf{Вопрос:} Верно ли, что все слова русского языка не содержат более 4 согласных подряд? Нет, «контрпример» имеет 5.

\subsection*{Прямое док-во}

Если $a \divs c$ и $b \divs d$, то $ab \divs cd$.
\[ \left.
\begin{aligned}
a \divs c &\implies \exists k: a = ck \\
b \divs d &\implies \exists l: b = dl
\end{aligned}
\right\} \implies ab = ckdl \implies ab \divs cd \text{ по опр.} \]

\subsection*{Док-во от противного}
\noindent \hfill \[ A \longrightarrow B \iff \overline{B} \longrightarrow \overline{A} \]
\noindent
\textbf{Пример:} Доказать, что простых чисел бесконечно много.
\begin{proof}[Док-во]
\noindent
$]$ не так, тогда $p_1, \dots, p_n$ --- все простые числа. \\
Р/м $x = p_1 \cdot p_2 \cdot p_3 \cdot \dots \cdot p_n + 1$, заметим, что 
$\left.
\begin{aligned}
x > p_i, \; \forall i \\
x \ndivs p_i, \; \forall i
\end{aligned}
\right\} \implies x \text{ --- новое простое число.}
\end{proof}

\subsection*{Оценка + пример}

\textbf{Вопрос:} Какое наибольшее кол-во не бьющих друг друга ладей?

\begin{proof}[Док-во]
Пример на 8:
\begin{center}
\begin{picture}(80,80)
\multiput(0,0)(0,10){9}{\line(1,0){80}}
\multiput(0,0)(10,0){9}{\line(0,1){80}}
\multiput(5,5)(10,10){8}{\circle*{4}}
\end{picture}
\end{center}

\noindent От противного: $] \, 9$, тогда 2 будут на одной горизонтали и будут бить друг друга. \\
\textit{На самом деле здесь мы пользуемся принципом Дирихле.}
\end{proof}

\subsection*{Принцип Дирихле}

\textbf{Утв.} Даны $n$ клеток и $n+1$ кроликов, то существует клетка с $\ge 2$ кроликами.

\paragraph{Утв. Обобщенный принцип Дирихле} 
Даны $n$ клеток, $k$ кроликов, то:
\begin{enumerate}
    \item $\exists$ клетка с $\ge \lceil \frac{k}{n} \rceil$ кроликами.
    \item $\exists$ клетка с $\le \lfloor \frac{k}{n} \rfloor$ кроликами.
\end{enumerate}

\begin{proof}[Док-во обобщенного принципа Дирихле:]
\noindent 
\begin{enumerate}
    \item От противного $]$ в каждой $< \lceil \frac{k}{n} \rceil$. В клетках $a_1, \dots, a_n$ кроликов. \\
    \[ a_i < \left\lceil \frac{k}{n} \right\rceil \implies a_i \le \left\lceil \frac{k}{n} \right\rceil - 1 < \frac{k}{n} \]
    \[ k = \sum a_i < n \cdot \frac{k}{n} = k \quad ?! \]

    \item $]$ в каждой $> \lfloor \frac{k}{n} \rfloor$.
    \[ a_i > \left \lfloor \frac{k}{n} \right \rfloor \implies a_i \ge \left \lfloor \frac{k}{n} \right \rfloor + 1 > \frac{k}{n} \]
    \[ k = \sum a_i > n \cdot \frac{k}{n} = k \quad ?! \]
\end{enumerate}
\end{proof}

\textbf{Пример:} 40 человек. Докажите, что есть месяц с $\ge 4$ днями рождения и есть месяц с $\le 3$. \\

\textit{Док-во по принципу Дирихле:}
\begin{enumerate}
    \item $]$ в каждом $\le 3 \implies$ всего дней рожд $\le 36$, но людей $40$?!
    
    \item $]$ в каждом $\ge 4 \implies$ всего $\ge 48$, но людей $40$?!
\end{enumerate}
\end{proof}

\paragraph{Утв. (Бесконечный принцип Дирихле)} 
Есть бесконечное число кроликов и конечное число клеток, тогда $\exists$ клетка, в которой бесконечное число кроликов.
\begin{proof}[Док-во:]
$ ] $ нет, тогда в каждой клетке конечное число кроликов, но $\text{конечное} \times \text{конечное} = \text{конечное} < \text{бесконечное}$.
\end{proof}

\subsection*{Комбинаторные док-ва}

Вместо счета объектов и проверки их равенства, можно установить соответствие.

\subsection*{Подсчет двумя способами}
Часто используется в док-ве от противного: «что-то одновременно больше и меньше нуля».

\subsection*{Принцип крайнего}

Хорошая штука, когда не хватает условий. Часто работает с методом от противного, но это не обязательно.
\noindent        
\textbf{Задача:} Докажите, что среди $n\ge 3$ точек на плоскости есть такая, что все углы с другими точками $<180^\circ$. 
\begin{proof}
    Рассмотрим самую левую точку и из таких самую нижнюю, заметим, что она подойдет, так как нет точек левее ее и 180 тоже не бывает. 
\end{proof}

\subsection*{Инвариант --- «что-то, что не меняется»}

\textbf{Задача:} На доске написано число $a$. Разрешается его стереть и написать либо $3a$, либо $a+8$, либо $\frac{a}{5}$, если $a \divs 5$. Можно ли из $1$ получить $2$?

\noindent \textit{Док-во:} Нет, т.к. число всегда нечётное. Инвариант --- чётность.

\noindent \textbf{Вопрос:} А из числа $2$ в $4$? Нет, но инвариант сложнее.

\subsection*{Вспомогательная раскраска}

\textbf{Задача:} Из шахматной доски вырезали 2 противоположных угла. Можно ли её разрезать на доминошки $1 \times 2$?

\noindent \textit{Док-во:} Отметим шахматную раскраску. Пусть отрезанные углы чёрные, тогда у нас 32 белых, 30 чёрных, а каждая доминошка содержит по одной клетке каждого цвета $\implies$ не выйдет.

\noindent Инвариант: кол-во белых $-\,\, \text{кол-во}$ чёрных.

\subsection*{Полуинвариант --- «что-то, что меняется лишь в одну сторону»}

\textbf{Пример:} предыдущая задача, но можно ещё вырезать чёрные клетки как действие.

\noindent Полуинвариант: кол-во белых $-\,\, \text{кол-во}$ чёрных. Он только увеличивается, изначально 2, а должен стать 0 ?!

\section*{Индукция}

Нужна, чтобы доказывать утверждение для всех $n$ ($\forall n$).\\

\noindent \textbf{Пример 1:} $1 + 2 + 3 + \dots + n = \frac{n(n+1)}{2}$

\noindent База: $n = 1 \implies 1 = \frac{1(1+1)}{2}$

\noindent Переход: $n = k \implies n = k+1$. Хотим для $k+1$:
\[ 1 + 2 + 3 + \dots + k + (k+1) = \underbrace{\frac{k(k+1)}{2}}_{\text{= по ИП}} + (k+1) = \frac{(k+1)(k+2)}{2} = \frac{(k+1)((k+1)+1)}{2} \quad \checkmark \]

\subsection*{Возвратная индукция --- переход из $1, \dots, n \implies n+1$}

\textbf{Пример:} Алгоритм Евклида заканчивается. Делаем индукцию по $k = n + m$.

\noindent Переход: $(n, m) = (n - m, m)$ --- для этой пары заканчивается по ИП (индукционному предположению), т.к. $(n - m) + m < n + m$.

\subsection*{Переход: $n \implies n+x$}

\textbf{Задача:} Какую сумму можно выдать монетами номинала 3 и 5? \\
Докажем, что любую сумму $n \ge 8$ можно выдать.

\noindent Переход: $n \implies n+3$. Выдадим 3, а $n$ --- по ИП.

\noindent База: 
$n = 8 = 3 + 5; \quad n = 9 = 3 + 3 + 3; \quad n = 10 = 5 + 5$.

\subsection*{Переход: $2^n \implies 2^{n+1}$ и $x \implies x-1$}

\textbf{Задача:} Доказать, что любое число раскладывать в сумму $2^k$.

\noindent Переход: $n \implies 2n$ и $n \implies 2n+1$.

\end{document}
